\clearpage
\fancyhead[L]{\small INHALT.}
\SourceLeaf{661}
\begingroup
\raggedright
\noindent\hangindent=1em Allgemeine Tafeln für Aberration und Nutation \mbox{.\ .\ .\ .\ .\ .\ .\ .\ .\ .} \hfill\mbox{1808 April}\quad\mbox{Seite 123}\par\smallskip
\noindent\hangindent=1em Ueber eine Aufgabe der sphaerischen Astronomie \mbox{.\ .\ .\ .\ .\ .\ .\ .\ .\ .} \hfill\mbox{1808 Octob.}\quad\mbox{--- 129}\par\smallskip
\noindent\hangindent=1em \textsc{Gauss} an von \textsc{Lindenau} [Zusatz zum vorhergehenden Aufsatz] \mbox{.\ .\ .\ .\ .} \hfill\mbox{1808 Novbr. 30}\quad\mbox{--- 141}\par\smallskip
\noindent\hangindent=1em \textsc{Gauss} an von \textsc{Lindenau} [Polhöhe-Bestimmung] \mbox{.\ .\ .\ .\ .\ .\ .\ .\ .\ .} \hfill\mbox{1808 Aug. 30}\quad\mbox{--- 146}\par\smallskip
\noindent\hangindent=1em Summarische Uebersicht der zur Bestimmung der Bahnen der beiden neuen Hauptplaneten angewandten Methoden \mbox{.\ .\ .\ .\ .\ .\ .\ .\ .\ .\ .\ .} \hfill\mbox{1809 Sept.}\quad\mbox{--- 148}\par\smallskip
\noindent\hangindent=1em Tafeln für die Mittags-Verbesserung \mbox{.\ .\ .\ .\ .\ .\ .\ .\ .\ .\ .\ .\ .\ .} \hfill\mbox{1811 März}\quad\mbox{--- 166}\par\smallskip
\noindent\hangindent=1em Tafel für die Sonnen-Coordinaten \mbox{.\ .\ .\ .\ .\ .\ .\ .\ .\ .\ .\ .\ .\ .\ .} \hfill\mbox{1812 Jan.}\quad\mbox{--- 172}\par\smallskip
\noindent\hangindent=1em Neue Aussicht zur Erweiterung des Gebiets der Himmelskunde \mbox{.\ .\ .\ .\ .} \hfill\mbox{1813}\quad\mbox{--- 181}\par\smallskip
\noindent\hangindent=1em Refractionstafeln \mbox{.\ .\ .\ .\ .\ .\ .\ .\ .\ .\ .\ .\ .\ .\ .\ .\ .\ .\ .\ .\ .\ .} \hfill\mbox{1822}\quad\mbox{--- 185}\par\smallskip
\noindent\hangindent=1em \textsc{Gauss} an \textsc{Schumacher} [Berechnung der wahren Anomalie eines Cometen] \mbox{.} \hfill\mbox{1843 April 1}\quad\mbox{--- 191}\par\smallskip
\noindent\hangindent=1em \textsc{Gauss} an \textsc{Schumacher} [Limiten eines Zodiacus] \mbox{.\ .\ .\ .\ .\ .\ .\ .\ .\ .} \hfill\mbox{1847 Nov. 23}\quad\mbox{--- 194}\par\smallskip
\par\medskip
\noindent\textit{Beobachtungen und Rechnungen.}\par\smallskip\nopagebreak[4]
\noindent\hangindent=1em Fortgesetzte Nachrichten über den längst vermutheten neuen Haupt-Planeten unseres Sonnen-Systems [Ceres Elemente I. II. III. IV.] \mbox{.\ .\ .\ .} \hfill\mbox{1801 Dec.}\quad\mbox{--- 199}\par\smallskip
\noindent\hangindent=1em Fortgesetzte Nachrichten über den neuen Haupt-Planeten [Ceres. Wiederauffindung] \mbox{.\ .\ .\ .\ .\ .\ .\ .\ .\ .\ .\ .\ .\ .\ .\ .\ .\ .\ .\ .\ .\ .} \hfill\mbox{1802 Febr.}\quad\mbox{--- 204}\par\smallskip
\noindent\hangindent=1em Fortgesetzte Nachrichten über den neuen Haupt-Planeten. Ceres Ferdinandea [Ceres Elemente V. VI. VII.] \mbox{.\ .\ .\ .\ .\ .\ .\ .\ .\ .\ .\ .\ .\ .} \hfill\mbox{1802 März}\quad\mbox{--- 205}\par\smallskip
\noindent\hangindent=1em Fortgesetzte Nachrichten. [Ceres. Beobachtungen] \mbox{.\ .\ .\ .\ .\ .\ .\ .\ .} \hfill\mbox{1802 April}\quad\mbox{--- 209}\par\smallskip
\noindent\hangindent=1em Fortgesetzte Nachrichten. [Ceres. Verglichene Beobachtungen] \mbox{.\ .\ .\ .} \hfill\mbox{1802 Mai}\quad\mbox{--- 211}\par\smallskip
\noindent\hangindent=1em Fortgesetzte Nachrichten. [Pallas. Elemente I. II.] \mbox{.\ .\ .\ .\ .\ .\ .\ .} \hfill\mbox{1802 Juni}\quad\mbox{--- 212}\par\smallskip
\noindent\hangindent=1em Fortgesetzte Nachrichten. [Pallas. Elemente III.] \mbox{.\ .\ .\ .\ .\ .\ .\ .\ .} \hfill\mbox{1802 Juli}\quad\mbox{--- 217}\par\smallskip
\noindent\hangindent=1em Pallas Olbersiana [Vergleichung mit den Elementen III.] \mbox{.\ .\ .\ .\ .\ .\ .} \hfill\mbox{1802 Aug.}\quad\mbox{--- 221}\par\smallskip
\noindent\hangindent=1em \textsc{Gauss} an \textsc{Bode}. [Pallas. Elemente IV.] \mbox{.\ .\ .\ ,\ .\ .\ .\ .\ .\ .\ .\ .\ .} \hfill\mbox{1802 Sept. 10}\quad\mbox{--- 223}\par\smallskip
\noindent\hangindent=1em Ceres Ferdinandea. [Beobachtungen und Vergleichungen] \mbox{.\ .\ .\ .\ .\ .} \hfill\mbox{1802 Juni}\quad\mbox{--- 224}\par\smallskip
\noindent\hangindent=1em Ceres Ferdinandea. [Störungen durch Jupiter. Länge und Radius Vector. Elemente VIII.] \mbox{.\ .\ .\ .\ .\ .\ .\ .\ .\ .\ .\ .\ .\ .\ .\ .\ .\ .\ .\ .} \hfill\mbox{1802 Octob.}\quad\mbox{--- 224}\par\smallskip
\noindent\hangindent=1em Pallas Olbersiana. [Ephemeride] \mbox{.\ .\ .\ .\ .\ .\ .\ .\ .\ .\ .\ .\ .\ .\ .\ .} \hfill\mbox{1802 Octob.}\quad\mbox{--- 226}\par\smallskip
\noindent\hangindent=1em Ceres Ferdinandea. [Störung durch Jupiter. Länge, Radius Vector und Breite] \hfill\mbox{1802 Nov.}\quad\mbox{--- 227}\par\smallskip
\noindent\hangindent=1em Pallas Olbersiana. [Abstände der Planeten von der Sonne] \mbox{.\ .\ .\ .\ .\ .} \hfill\mbox{1802 Nov.}\quad\mbox{--- 230}\par\smallskip
\noindent\hangindent=1em Vorübergang des Mercur vor der Sonne, den 3. November 1802 \mbox{.\ .\ .\ .\ .} \hfill\mbox{1802 Dec.}\quad\mbox{--- 231}\par\smallskip
\noindent\hangindent=1em Pallas Olbersiana. [Elemente V. und Vergleichungen] \mbox{.\ .\ .\ .\ .\ .\ .\ .} \hfill\mbox{1802 Dec.}\quad\mbox{--- 232}\par\smallskip
\noindent\hangindent=1em \textsc{Gauss} an \textsc{Bode}. [Pallas. Elemente VI.] \mbox{.\ .\ .\ .\ .\ .\ .\ .\ .\ .\ .\ .\ .} \hfill\mbox{1803 März 3}\quad\mbox{--- 234}\par\smallskip
\noindent\hangindent=1em Tafeln für die Störungen der Ceres durch Jupiter \mbox{.\ .\ .\ .\ .\ .\ .\ .\ .\ .} \hfill\mbox{1802 Dec. 26}\quad\mbox{--- 235}\par\smallskip
\noindent\hangindent=1em Pallas. [Beobachtungen] \mbox{.\ .\ .\ .\ .\ .\ .\ .\ .\ .\ .\ .\ .\ .\ .\ .\ .\ .} \hfill\mbox{1803 April u. Mai}\quad\mbox{--- 244}\par\smallskip
\noindent\hangindent=1em Pallas. [Beobachtungen] \mbox{.\ .\ .\ .\ .\ .\ .\ .\ .\ .\ .\ .\ .\ .\ .\ .\ .\ .\ .} \hfill\mbox{1803 Juni}\quad\mbox{'--- 244}\par\smallskip
\noindent\hangindent=1em Pallas. [Beobachtungen verglichen mit Elementen VI.] \mbox{.\ .\ .\ .\ .\ .\ .} \hfill\mbox{1803 Juli}\quad\mbox{--- 245}\par\smallskip
\noindent\hangindent=1em Ceres. [Elemente IX.] \mbox{.\ .\ .\ .\ .\ .\ .\ .\ .\ .\ .\ .\ .\ .\ .\ .\ .\ .\ .} \hfill\mbox{1803 Sept. u. Oct.}\quad\mbox{--- 245}\par\smallskip
\noindent\hangindent=1em Ceres und Pallas. [Pallas Elemente VII.] \mbox{.\ .\ .\ .\ .\ .\ .\ .\ .\ .\ .\ .} \hfill\mbox{1804 März}\quad\mbox{--- 246}\par\smallskip
\noindent\hangindent=1em Neuer Comet [von 1804. Elemente] \mbox{.\ .\ ,\ .\ .\ .\ .\ .\ .\ .\ .\ .\ .\ .\ .} \hfill\mbox{1804 Mai}\quad\mbox{--- 247}\par\smallskip
\noindent\hangindent=1em \textsc{Gauss} an \textsc{Bode}. [Juno. Elemente I. III.] \mbox{.\ .\ .\ .\ .\ .\ .\ .\ .\ .} \hfill\mbox{1804 Sept. 30 u. Nov. 12}\quad\mbox{--- 248}\par\smallskip
\noindent\hangindent=1em Ueber einen Wandelstern. [Juno. Elemente I.) \mbox{.\ .\ ,\ .\ .\ .\ .\ .\ .\ .} \hfill\mbox{1804 Octob.}\quad\mbox{--- 250}\par\smallskip
\noindent\hangindent=1em Einige Nachrichten über den neuen Planeten. [Juno] \mbox{.\ .\ ,\ .\ ,\ .\ .\ .} \hfill\mbox{1804 Octob. 6}\quad\mbox{--- 253}\par\smallskip
\endgroup
